\chapter{RESULTADOS ESPERADOS}

Com o desenvolvimento deste trabalho, pretende-se construir um cenário de simulação microscópica baseado em um recorte urbano real da região central de Videira, permitindo analisar, de forma controlada, o impacto de diferentes perfis de motoristas sobre o desempenho do tráfego.

Como resultado inicial, busca-se obter uma rede viária funcional no \gls{sumo}, configurada de forma simplificada, mas suficiente para representar a dinâmica principal de circulação da área estudada. Também deverão ser definidos perfis de motoristas com características distintas de condução, possibilitando a criação de cenários comparativos com diferentes níveis de fluxo e diferentes composições de comportamento dos condutores.

Do ponto de vista analítico, o trabalho deverá permitir a comparação entre os perfis cauteloso, normal e agressivo a partir de métricas como tempo médio de viagem, velocidade média, tempo parado, formação de filas e vazão de veículos. A análise dos cenários simulados poderá indicar se determinados comportamentos produzem efeitos mais significativos sobre o tráfego, especialmente em situações de maior carregamento da rede.

Além disso, os resultados poderão contribuir para compreender como a variação de parâmetros comportamentais influencia o desempenho de um sistema viário urbano simulado. Dessa forma, a contribuição do trabalho não está em validar o uso do \gls{sumo} como ferramenta, mas em aplicar a simulação computacional para investigar um problema específico relacionado ao comportamento dos motoristas em um cenário urbano delimitado.

Por fim, pretende-se que a pesquisa forneça uma base metodológica para a continuidade do estudo no TCC, incluindo a configuração dos perfis, a execução dos cenários, a coleta das métricas e a análise comparativa dos resultados. Essa base poderá ser ampliada futuramente com novos parâmetros, outros recortes urbanos, diferentes tipos de veículos ou condições adicionais de simulação.
